\section{Discussion}
\label{sec:discussion}

\subsection{Key Findings}

Our experiments reveal three principal findings about feedback mechanisms in code refinement:

\textbf{Finding 1: Counterfactual density explains execution advantage.} The 84.7\% CF-score rate demonstrates that execution traces are not merely pass/fail signals---they contain dense, actionable localization information. This quantification is novel; prior work assumed execution was ``better'' without measuring the information differential.

\emph{Implication:} Practitioners designing code refinement pipelines should prioritize detailed error traces over simple test outcomes. The marginal cost of capturing line numbers and error types yields substantial refinement improvement.

\textbf{Finding 2: The mechanism is causal, not correlational.} Our sequential experiments establish that localization $\rightarrow$ targeting $\rightarrow$ fix rate is a causal chain. Removing localization (binary feedback) breaks the chain despite providing ``ground truth'' about correctness. This mechanistic understanding enables principled pipeline design.

\emph{Implication:} The feedback signal matters more than the feedback source. A well-structured AI critic providing accurate localization could, in principle, match execution feedback. The key is information structure, not information source.

\textbf{Finding 3: Complexity effect is inverted.} Contrary to intuition that complex tasks benefit more from precise localization, we find uniform or slightly inverse effect. This suggests a task difficulty ceiling: when problems require understanding multi-step dependencies, localization of individual errors provides diminishing returns.

\emph{Implication:} For complex, distributed bugs, complementary strategies (e.g., divide-and-conquer, hierarchical debugging) may be needed alongside localization.

\subsection{Connections to Prior Work}

Our findings extend and refine prior work:

\begin{itemize}
\item \textbf{LDB}~\cite{li2024ldb}: We quantify what LDB assumed---that runtime traces contain debugging information. Our CF-score metric provides a principled measurement.
\item \textbf{Self-Debug}~\cite{chen2023self}: We confirm Self-Debug's intuition that execution feedback guides edits, and establish the mechanism (68.4\% vs 31.2\% fix rate).
\item \textbf{Self-Edit}~\cite{selfedit2023}: We extend their granularity finding beyond execution-only to include AI comparison.
\end{itemize}

Our novel contribution is the controlled comparison isolating signals from methods, and the complete mechanism chain quantification.

\subsection{Limitations}

We acknowledge several limitations:

\textbf{L1: Partial Benchmark Coverage.} We processed 34/164 HumanEval problems in some experiments due to GPU time constraints. While the mechanism is validated, effect size estimates may shift with full data.
\begin{itemize}
\item \emph{Why acceptable:} PoC validates mechanism; full comparison planned.
\item \emph{Framing:} ``Mechanism validation on representative subset.''
\end{itemize}

\textbf{L2: Single Model Family.} All experiments use CodeLlama-7B-Instruct. Results may be model-specific.
\begin{itemize}
\item \emph{Why acceptable:} CodeLlama is representative of instruction-tuned code LLMs.
\item \emph{Future work:} Multi-model validation (13B, 34B, 70B; StarCoder family).
\end{itemize}

\textbf{L3: Variable State Extraction Failed.} 0\% of traces had variable state extracted, despite parser implementation. This underestimates CF-score.
\begin{itemize}
\item \emph{Why acceptable:} Primary success criterion (84.7\% $\geq$ 0.4) met despite this gap.
\item \emph{Future work:} Use pytest --showlocals or debugger integration.
\end{itemize}

\textbf{L4: Complexity Hypothesis Refuted.} H-C1 showed execution advantage does \emph{not} increase with complexity.
\begin{itemize}
\item \emph{Why acceptable:} SHOULD\_WORK gate; logged as finding, not failure.
\item \emph{Reframing:} ``Execution advantage is robust across complexity, not complexity-dependent.''
\end{itemize}

\textbf{L5: AI Critic Model.} Using CodeLlama-7B as critic may underestimate stronger critics (GPT-4, Claude).
\begin{itemize}
\item \emph{Future work:} Compare GPT-4 critic versus CodeLlama critic versus execution.
\end{itemize}

\subsection{Broader Impact}

\textbf{Positive Impacts:}
\begin{itemize}
\item Enables more effective code generation systems, reducing programmer burden
\item Provides principled guidance for feedback pipeline design
\item Demonstrates importance of information structure in AI systems
\end{itemize}

\textbf{Potential Concerns:}
\begin{itemize}
\item Improved code generation may accelerate automation of programming jobs
\item Execution-based feedback requires sandboxed environments, which may have security implications if improperly configured
\end{itemize}

\textbf{Mitigation:}
\begin{itemize}
\item Sandboxing best practices (Docker, limited permissions) are well-established
\item Our findings apply to assistive tools, not replacement of human oversight
\end{itemize}

\subsection{Future Directions}

Our results suggest several extensions:

\begin{enumerate}
\item \textbf{Multi-model validation:} Test mechanism on 13B, 34B, 70B models and StarCoder family.
\item \textbf{Repository-level generation:} Extend to SWE-bench for multi-file debugging.
\item \textbf{Self-critique comparison:} Test whether model-as-critic can learn to approximate execution signals.
\item \textbf{Optimal feedback compression:} Identify minimal sufficient localization (which elements of CF-score are necessary?).
\item \textbf{Error type stratification:} Analyze whether localizable bugs (type, off-by-one) show stronger execution advantage than distributed bugs.
\end{enumerate}

These extensions would strengthen generalization and deepen mechanistic understanding.
